\documentclass[12pt]{article}

\usepackage{amsmath,amssymb,amsthm}
\usepackage{geometry}
\usepackage{setspace}
\usepackage{enumitem}
\usepackage{hyperref}

\geometry{letterpaper, margin=1in}
\setstretch{1.15}

\newtheorem{theorem}{Theorem}[section]
\newtheorem{proposition}[theorem]{Proposition}
\newtheorem{corollary}[theorem]{Corollary}
\newtheorem{lemma}[theorem]{Lemma}

\theoremstyle{definition}
\newtheorem{definition}[theorem]{Definition}

\theoremstyle{remark}
\newtheorem{remark}[theorem]{Remark}
\newtheorem{observation}[theorem]{Observation}

\title{\textbf{Three Independent Paths to a Single Constant: \\[0.3cm]
Period Descent and the Convergence of \\
Lemniscatic, Modular, and Lattice Invariants}}

\author{W.\,J.\ cpaci-tani~III}
\date{}

\begin{document}

\maketitle

\begin{abstract}
\noindent We identify three mathematical objects---the lemniscate constant $\varpi$, the Jacobi theta function $\theta_3$ at the self-dual point, and the Watson integral $W_3$ of the three-dimensional cubic lattice---each arising from a different branch of mathematics with no reference to the others. We prove that all three evaluate to expressions in $\Gamma(1/4)$ via methodologically independent derivation paths, and show that the convergence has a structural explanation: all three are period integrals of the CM elliptic curve $E\!: y^2 = x^3 - x$. We further prove a Diophantine constraint on when the transcendental period $\pi$ can be eliminated from a mixed continuous-discrete trace, establishing that the fundamental solution $(p{=}3,\,q{=}2)$ forces the result into the purely lemniscatic period ring. No physical assumptions are required. Every result is a theorem of pure mathematics.
\end{abstract}

\tableofcontents

%------------------------------------------------------------------
\section{Introduction}
%------------------------------------------------------------------

A recurring theme in analytic number theory is the appearance of the Euler gamma value $\Gamma(1/4)$ in seemingly unrelated contexts: elliptic integrals, modular forms, and lattice Green's functions. We show that these appearances are not coincidental. They are controlled by a common geometric ancestor---the CM elliptic curve $E\!: y^2 = x^3 - x$ with endomorphism ring $\mathbb{Z}[i]$---and by a Diophantine condition on dimensions that determines when the transcendental period $\pi$ can be eliminated from mixed continuous-discrete traces.

Our exposition is deliberately self-contained with respect to the algebraic manipulations. We define three objects from three independent mathematical domains, show they produce the same numerical constant, identify their common origin, and prove the dimension constraint. Where we invoke deep classical results (Watson's 1939 evaluation, the Chowla--Selberg formula, Chudnovsky's algebraic independence theorem), we state them precisely and give references. No physics, no lattice field theory, and no prior knowledge of any particular theoretical framework is assumed or required.

%------------------------------------------------------------------
\section{Three Independent Definitions}
\label{sec:three}
%------------------------------------------------------------------

We define three mathematical quantities. Each definition uses only the standard apparatus of its parent field. No definition references either of the other two.

%------------------------------------------------------------------
\subsection{Path A: The Lemniscate Integral (Algebraic Geometry)}

\begin{definition}[The Lemniscatic Constant]
\label{def:varpi}
The \emph{lemniscate constant} $\varpi$ is twice the arc length of one lobe of the Bernoulli lemniscate $r^2 = \cos 2\theta$:
\begin{equation}
    \varpi \;=\; 2\!\int_0^1 \frac{dx}{\sqrt{1-x^4}}
\end{equation}
Equivalently, $\varpi = 2I_4$ where $I_4 = \int_0^1 (1-x^4)^{-1/2}\,dx$ is the lemniscate quarter-arc.
\end{definition}

This integral is a special value of the Beta function. Substituting $u = x^4$:
\begin{equation}
    I_4 = \tfrac{1}{4}B\!\left(\tfrac{1}{4},\,\tfrac{1}{2}\right) = \frac{\Gamma(1/4)\,\Gamma(1/2)}{4\,\Gamma(3/4)}
\end{equation}
We apply the Euler reflection formula $\Gamma(z)\Gamma(1-z) = \pi/\sin(\pi z)$. At $z = 1/4$:
\begin{equation}
    \Gamma(1/4)\,\Gamma(3/4) = \frac{\pi}{\sin(\pi/4)} = \frac{\pi}{1/\sqrt{2}} = \pi\sqrt{2}
\end{equation}
Substituting $\Gamma(3/4) = \pi\sqrt{2}/\Gamma(1/4)$ and $\Gamma(1/2) = \sqrt{\pi}$ into the expression for $I_4$:
\begin{equation}
    I_4 = \frac{\Gamma(1/4)\cdot\sqrt{\pi}}{4}\cdot\frac{\Gamma(1/4)}{\pi\sqrt{2}} = \frac{\Gamma(1/4)^2}{4\sqrt{2\pi}}
\end{equation}
Therefore:
\begin{equation}
    \boxed{\;\varpi \;=\; 2I_4 \;=\; \frac{\Gamma(1/4)^2}{2\sqrt{2\pi}}\;}
    \label{eq:varpi}
\end{equation}

\begin{remark}
The \emph{only} inputs are: the definition of the lemniscate curve, the Beta integral evaluation, and the Euler reflection formula. No theta functions, no lattices, no modular forms.
\end{remark}

%------------------------------------------------------------------
\subsection{Path B: The Jacobi Theta Function (Modular Forms)}

\begin{definition}[Theta Function at Self-Dual Point]
\label{def:theta}
The Jacobi theta function is
\begin{equation}
    \theta_3(q) = \sum_{n=-\infty}^{\infty} q^{n^2}
\end{equation}
We evaluate at $q = e^{-\pi}$, the self-dual point under the Poisson summation $\tau \mapsto -1/\tau$ (equivalently, $\tau = i$).
\end{definition}

The Chowla--Selberg formula \cite{chowla-selberg} evaluates theta functions of CM lattices in terms of gamma values. At $\tau = i$, the lattice $\mathbb{Z} + i\mathbb{Z}$ is the period lattice of the CM curve $E\!: y^2 = x^3 - x$ with endomorphism ring $\mathbb{Z}[i]$. The specialization gives:
\begin{equation}
    \theta_3(e^{-\pi}) = \frac{\pi^{1/4}}{\Gamma(3/4)}
    \label{eq:theta-raw}
\end{equation}

Applying the Euler reflection formula as in Path A, $\Gamma(3/4) = \pi\sqrt{2}/\Gamma(1/4)$, so:
\begin{equation}
    \boxed{\;\theta_3(e^{-\pi}) \;=\; \frac{\Gamma(1/4)}{\sqrt{2}\;\pi^{3/4}}\;}
    \label{eq:theta}
\end{equation}

\begin{remark}
The \emph{only} inputs specific to this path are: the definition of $\theta_3$ and the Chowla--Selberg formula at $\tau = i$. The Euler reflection formula is shared with Path A, but the objects being evaluated ($\theta_3$ vs.\ the arc length integral) are distinct. No lemniscate geometry or lattice Green's functions are used.
\end{remark}

%------------------------------------------------------------------
\subsection{Path C: The Watson Integral (Lattice Combinatorics)}

\begin{definition}[Watson's Integral for the Cubic Lattice]
\label{def:watson}
The \emph{Watson integral} $W_3$ is the lattice Green's function of the simple cubic lattice $\mathbb{Z}^3$ evaluated at the origin:
\begin{equation}
    W_3 = \frac{1}{(2\pi)^3}\!\int_{[-\pi,\pi]^3}\!\!\frac{d^3k}{3 - \cos k_1 - \cos k_2 - \cos k_3}
\end{equation}
\end{definition}

Watson \cite{watson1939} evaluated this integral by iterated contour integration, reducing the Brillouin zone integral to a product of complete elliptic integrals. The complete elliptic integral of the first kind at the lemniscatic modulus $k = 1/\sqrt{2}$ satisfies $K(1/\sqrt{2}) = \Gamma(1/4)^2/(4\sqrt{\pi})$ (a classical identity due to Legendre). Watson's reduction yields:
\begin{equation}
    \boxed{\;W_3 \;=\; \frac{\Gamma(1/4)^4}{4\pi^3}\;}
    \label{eq:watson}
\end{equation}

\begin{remark}
The \emph{only} inputs are: the definition of the lattice Laplacian on $\mathbb{Z}^3$, Watson's 1939 contour integration method, and the Legendre evaluation of $K(1/\sqrt{2})$. No lemniscate arc lengths, no theta functions, no modular forms.
\end{remark}

%------------------------------------------------------------------
\section{The Convergence}
\label{sec:convergence}
%------------------------------------------------------------------

We now show that the three independently defined quantities are algebraically related.

\begin{proposition}[Algebraic Relations]
\label{prop:verified}
Define $\Phi = \Gamma(1/4)^2$. Then:
\begin{align}
    \varpi &= \frac{\Phi}{2\sqrt{2\pi}} \label{eq:v1}\\[4pt]
    \theta_3(e^{-\pi})^2 &= \frac{\Phi}{2\pi^{3/2}} \label{eq:v2}\\[4pt]
    W_3 &= \frac{\Phi^2}{4\pi^3} \label{eq:v3}
\end{align}
In particular, eliminating $\Phi$:
\begin{align}
    \theta_3(e^{-\pi})^2 &= \frac{\sqrt{2}\,\varpi}{\pi} \label{eq:cross1}\\[4pt]
    W_3 &= \frac{2\varpi^2}{\pi^2} \label{eq:cross2}\\[4pt]
    W_3 &= \theta_3(e^{-\pi})^4 \label{eq:cross3}
\end{align}
\end{proposition}

\begin{proof}
Equations~\eqref{eq:v1}--\eqref{eq:v3} are the evaluations from Section~\ref{sec:three}.

\emph{Equation~\eqref{eq:cross1}:} From~\eqref{eq:v1}, $\Phi = 2\sqrt{2\pi}\,\varpi$. Substituting into~\eqref{eq:v2}:
\[
    \theta_3^2 = \frac{2\sqrt{2\pi}\,\varpi}{2\pi^{3/2}} = \frac{\sqrt{2}\,\varpi}{\pi}. \qquad\checkmark
\]

\emph{Equation~\eqref{eq:cross2}:} Substituting $\Phi = 2\sqrt{2\pi}\,\varpi$ into~\eqref{eq:v3}:
\[
    W_3 = \frac{(2\sqrt{2\pi}\,\varpi)^2}{4\pi^3} = \frac{8\pi\,\varpi^2}{4\pi^3} = \frac{2\varpi^2}{\pi^2}. \qquad\checkmark
\]

\emph{Equation~\eqref{eq:cross3}:} From~\eqref{eq:v2}, $\Phi = 2\pi^{3/2}\,\theta_3^2$. Substituting into~\eqref{eq:v3}:
\[
    W_3 = \frac{(2\pi^{3/2}\,\theta_3^2)^2}{4\pi^3} = \frac{4\pi^3\,\theta_3^4}{4\pi^3} = \theta_3(e^{-\pi})^4. \qquad\checkmark\qedhere
\]
\end{proof}

\begin{remark}[Non-circularity]
\label{rmk:noncircular}
The critical point: \emph{each path arrives at $\Gamma(1/4)^2$ through entirely different mathematical machinery.}
\begin{itemize}[nosep]
    \item Path A uses the arc length of an algebraic curve and the Beta integral.
    \item Path B uses the heat kernel of a flat torus and the Chowla--Selberg formula.
    \item Path C uses the Green's function of the cubic lattice and Watson's contour integration.
\end{itemize}
Paths A and B share the Euler reflection formula as a simplification tool, but the objects being evaluated are distinct (an arc-length integral vs.\ a lattice sum). Path C is methodologically independent of both. The common output $\Gamma(1/4)^2$ is independently derived in each case. The question this paper addresses is: \emph{what structural fact explains the convergence?}
\end{remark}

%------------------------------------------------------------------
\section{The Structural Explanation: A Common CM Curve}
\label{sec:cm}
%------------------------------------------------------------------

The convergence of Paths A, B, and C is explained by the following:

\begin{theorem}[Common Period Origin]
\label{thm:cm}
Each of the three quantities---$\varpi$, $\theta_3(e^{-\pi})$, and $W_3$---is a period integral of the CM elliptic curve $E\!: y^2 = x^3 - x$.
\end{theorem}

The curve $E$ has the following distinguished properties:
\begin{itemize}[nosep]
    \item Complex multiplication by $\mathbb{Z}[i]$ (Gaussian integers)
    \item $j$-invariant $j(E) = 1728 = 12^3$
    \item Automorphism group $\operatorname{Aut}(E) \cong \mathbb{Z}/4\mathbb{Z}$ (maximal for elliptic curves over $\mathbb{Q}$)
    \item Conductor $N = 32$, minimal discriminant $\Delta = -64$
\end{itemize}

\begin{proof}[Proof sketch]
The real period of $E$ is $\Omega_1 = 2\!\int_0^1 dx/\sqrt{1-x^4} = \varpi$. This is Path A directly.

For Path B: the period lattice of $E$ is $\Lambda = \Omega_1(\mathbb{Z} + i\mathbb{Z})$, with $\tau = i$. The theta function of this lattice is $\theta_3(q)$ at $q = e^{-\pi|\tau|} = e^{-\pi}$, and the Chowla--Selberg formula evaluates theta functions of CM lattices in terms of gamma values at rational arguments determined by the CM discriminant.

For Path C: Watson's iterated contour integration reduces $W_3$ to products of $K(k)$ at specific moduli. The key modulus is $k = 1/\sqrt{2}$, the lemniscatic modulus, for which:
\[
    K(1/\sqrt{2}) = \frac{\Gamma(1/4)^2}{4\sqrt{\pi}} = \frac{\varpi\sqrt{2}}{2}
\]
(verification: $\varpi\sqrt{2}/2 = \Gamma(1/4)^2\sqrt{2}/(2\cdot 2\sqrt{2\pi}) = \Gamma(1/4)^2/(4\sqrt{\pi})$\;\checkmark). Therefore $W_3$ is an algebraic expression in the periods of $E$.
\end{proof}

The content of Theorem~\ref{thm:cm} is that these three a priori unrelated evaluation results all produce periods of a single curve. This is the structural explanation for the convergence: there is no coincidence to explain, because there is only one underlying object.

%------------------------------------------------------------------
\section{The Period Descent Condition}
\label{sec:descent}
%------------------------------------------------------------------

We now prove a result that characterizes when the transcendental period $\pi$ can be eliminated from mixed continuous-discrete traces, yielding purely lemniscatic (i.e., CM-period) expressions.

\begin{definition}[The Mixed Gaussian Trace]
\label{def:mgt}
For positive integers $p$ and $q$, define:
\begin{equation}
    \mathcal{Z}(p,q) \;=\; \underbrace{\left(\int_{\mathbb{R}} e^{-x^2/2}\,dx\right)^{\!p}}_{\text{Gaussian volume of }\mathbb{R}^p} \;\times\; \underbrace{\left(\sum_{n\in\mathbb{Z}} e^{-\pi n^2}\right)^{\!q}}_{\text{Heat trace on }\mathbb{T}^q\text{ at }\tau=i}
\end{equation}
\end{definition}

The continuous factor evaluates to $(2\pi)^{p/2}$. The discrete factor evaluates via the Chowla--Selberg formula~\eqref{eq:theta} to $\theta_3(e^{-\pi})^q = [\Gamma(1/4)/(\sqrt{2}\,\pi^{3/4})]^q$. Therefore:

\begin{equation}
    \mathcal{Z}(p,q) = 2^{(p-q)/2}\;\pi^{(2p-3q)/4}\;\Gamma(1/4)^q
    \label{eq:Zpq}
\end{equation}

\begin{theorem}[Period Descent Condition]
\label{thm:descent}
Let $\pi$ and $\Gamma(1/4)$ be algebraically independent over $\overline{\mathbb{Q}}$ \emph{(Chudnovsky~\cite{chudnovsky})}. Then $\mathcal{Z}(p,q)$ lies in the ring $\overline{\mathbb{Q}}[\Gamma(1/4)]$ \emph{(}i.e., is free of $\pi$\emph{)} if and only if:
\begin{equation}
    \boxed{\;2p = 3q\;}
\end{equation}
\end{theorem}

\begin{proof}
\emph{(Sufficiency.)} If $2p = 3q$, then $(2p-3q)/4 = 0$ and~\eqref{eq:Zpq} gives $\mathcal{Z}(p,q) = 2^{(p-q)/2}\,\Gamma(1/4)^q$, which lies in $\mathbb{Q}(\sqrt{2})[\Gamma(1/4)]$ and is free of $\pi$.

\emph{(Necessity.)} Suppose $2p - 3q \neq 0$. Then~\eqref{eq:Zpq} expresses $\mathcal{Z}(p,q)$ as $c\cdot\pi^{(2p-3q)/4}\cdot\Gamma(1/4)^q$ for some $c \in \mathbb{Q}(\sqrt{2})$. If $\mathcal{Z}(p,q)$ were in $\overline{\mathbb{Q}}[\Gamma(1/4)]$, we could write $\pi^{(2p-3q)/4} = P(\Gamma(1/4))$ for some algebraic function $P$. This would constitute an algebraic relation between $\pi$ and $\Gamma(1/4)$ over $\overline{\mathbb{Q}}$, contradicting Chudnovsky's theorem.
\end{proof}

\begin{corollary}[The $(3,2)$ Trace]
The minimal positive integer solution to $2p = 3q$ is $(p,q) = (3,2)$. At this point:
\begin{equation}
    \mathcal{Z}(3,2) = \sqrt{2}\;\Gamma(1/4)^2 \;\in\; \mathbb{Q}(\sqrt{2})[\Gamma(1/4)]
\end{equation}
\end{corollary}

\begin{remark}[Dimensional Mismatch]
\label{rmk:mismatch}
The product $\sqrt{2\pi}\cdot\theta_3(e^{-\pi})^2$---which equals $\Gamma(1/4)^2/(\sqrt{2}\,\pi)$---corresponds to $\mathcal{Z}(1,2)$, with $\pi$-exponent $(2-6)/4 = -1$. Because $(1,2)$ fails the descent condition $2(1) \neq 3(2)$, the factor of $\pi$ in the denominator cannot be eliminated. By contrast, $\mathcal{Z}(3,2) = \sqrt{2}\,\Gamma(1/4)^2$ is $\pi$-free, illustrating that the same period (of the same CM curve) can appear with or without $\pi$ depending on the dimensional context.
\end{remark}

%------------------------------------------------------------------
\section{Representations and the Bridge Constant}
\label{sec:representations}
%------------------------------------------------------------------

For applications in other work, it is convenient to define a single constant that encodes the relationship between the lemniscatic period and the circular period:

\begin{definition}[The Bridge Constant]
\begin{equation}
    G^* \;=\; \frac{2\varpi}{\sqrt{\pi}} \;=\; \frac{\sqrt{2}\;\Gamma(1/4)^2}{2\pi} \;\approx\; 2.9586751\ldots
    \label{eq:gstar}
\end{equation}
\end{definition}

The normalization $2\varpi/\sqrt{\pi}$ is not arbitrary: it is the unique rescaling of $\varpi$ such that $G^{*2} = 2\pi W_3$, directly equating the bridge constant to the Watson integral. This can be verified:
\[
    \frac{G^{*2}}{2\pi} = \frac{4\varpi^2}{\pi}\cdot\frac{1}{2\pi} = \frac{2\varpi^2}{\pi^2} = W_3 \qquad\text{(by eq.~\eqref{eq:cross2})}.
\]

We collect all known representations of $G^*$:

\begin{center}
\renewcommand{\arraystretch}{1.5}
\begin{tabular}{lll}
\hline
\textbf{Representation} & \textbf{Source Domain} & \textbf{Key Input} \\
\hline
$G^* = \dfrac{2\varpi}{\sqrt{\pi}}$ & Algebraic geometry & Lemniscate arc length \\[6pt]
$G^* = 2\sqrt{\pi}\;/\;\operatorname{agm}(1,\sqrt{2})$ & Arithmetic & Gauss's constant \\[6pt]
$G^* = \sqrt{2\pi}\;\theta_3(e^{-\pi})^2$ & Modular forms & Chowla--Selberg \\[6pt]
$G^* = \dfrac{1}{\sqrt{2\pi}}\,B(1/4,\,1/4)$ & Special functions & Beta function \\[6pt]
$G^* = \dfrac{2\sqrt{2}}{\sqrt{\pi}}\,K(1/\sqrt{2})$ & Elliptic integrals & Complete $K$ \\[6pt]
$G^{*2} = 2\pi\,W_3$ & Lattice combinatorics & Watson integral \\[6pt]
\hline
\end{tabular}
\end{center}

\begin{remark}
Each row can be verified independently using only the tools of its source domain. The convergence is a consequence of Theorem~\ref{thm:cm}: all these domains probe different aspects of the CM curve $E$.
\end{remark}

%------------------------------------------------------------------
\section{Summary}
%------------------------------------------------------------------

\begin{enumerate}
    \item Three independently defined mathematical objects---the lemniscate constant, the theta function at self-duality, and the Watson integral of $\mathbb{Z}^3$---all evaluate to expressions in $\Gamma(1/4)^2$.

    \item The structural explanation is Theorem~\ref{thm:cm}: all three are period integrals of the CM elliptic curve $E\!: y^2 = x^3 - x$, the unique elliptic curve over $\mathbb{Q}$ with $j = 1728$ and maximal automorphism group.

    \item The Period Descent Condition (Theorem~\ref{thm:descent}) characterizes when $\pi$ can be eliminated from mixed continuous-discrete traces: precisely when $2p = 3q$. The fundamental solution $(3,2)$ yields a $\pi$-free invariant in the period ring of $E$.

    \item No physics is assumed. Every result relies on: the Beta integral, the Chowla--Selberg formula, Chudnovsky's algebraic independence theorem, Watson's 1939 evaluation, and standard properties of elliptic curves with complex multiplication.
\end{enumerate}

%------------------------------------------------------------------
\begin{thebibliography}{9}
%------------------------------------------------------------------

\bibitem{watson1939}
G.\,N.\ Watson, ``Three triple integrals,'' \textit{Quart.\ J.\ Math.\ Oxford} \textbf{10} (1939), 266--276.

\bibitem{chudnovsky}
G.\,V.\ Chudnovsky, ``Contributions to the theory of transcendental numbers,'' \textit{Math.\ Surveys and Monographs} \textbf{19}, AMS (1984).

\bibitem{chowla-selberg}
S.\ Chowla and A.\ Selberg, ``On Epstein's zeta function,'' \textit{J.\ Reine Angew.\ Math.} \textbf{227} (1967), 86--110.

\bibitem{kontsevich-zagier}
M.\ Kontsevich and D.\ Zagier, ``Periods,'' in \textit{Mathematics Unlimited---2001 and Beyond}, Springer (2001), 771--808.

\bibitem{silverman}
J.\,H.\ Silverman, \textit{The Arithmetic of Elliptic Curves}, 2nd ed., Springer GTM~\textbf{106} (2009).

\bibitem{borwein-bailey}
J.\,M.\ Borwein and D.\,H.\ Bailey, \textit{Mathematics by Experiment}, 2nd ed., A K Peters (2008).

\end{thebibliography}

\end{document}
